\documentclass{article}

\usepackage{graphicx}
\graphicspath{ {../imgs/} }

\usepackage{makecell}
\usepackage[margin=0.5in]{geometry}
\usepackage{amssymb}
\usepackage{hyperref}
\usepackage[absolute,overlay]{textpos}
\usepackage[italian]{babel}

% Rimuovo l'indentazione automatica dei paragrafi
\usepackage{parskip}

\title{\textbf{Gara di Informatica \\~\\ Regolamento di Gara} \\~\\
\large DISIM - Università degli Studi dell'Aquila \\~\\
\large Street Science 2025 \\~\\
}

\date{}


\begin{document}
\begin{textblock*}{3cm}(1cm,1.75cm)
   \includegraphics[width=3cm]{imgs/univaq.png}
\end{textblock*}

\begin{textblock*}{4cm}(16cm,2cm)
   \includegraphics[width=3.5cm]{imgs/streetscience.jpg}
\end{textblock*}

\maketitle

\section*{}

\begin{enumerate}
    \item Ogni squadra è composta di massimo \textbf{6 componenti}: massimo \textbf{4 titolari}, che partecipano effettivamente alla gara, e \textbf{2 riserve}. I componenti di una squadra vengono preventivamente riconosciuti mediante un documento di identità, per consentire le verifiche di correttezza, e l'accesso all'aula di gara è consentito ai soli titolari
    \item Ciascuna squadra affronta la gara sul proprio \textbf{\textit{PC di gara}}, che può essere un dispositivo portato da casa o uno fornito dagli organizzatori. Nel primo caso, la squadra deve \textbf{segnalarlo} al team di Street Science
    \item La \textbf{connessione} in rete del \textit{PC di gara} può avvenire sfruttando l'hotspot personale di un proprio dispositivo (anche aggiuntivo, a patto che venga \textbf{segnalato} agli organizzatori) oppure, su richiesta, utilizzando la rete di Ateneo UnivAQ. In quest'ultimo caso, i partecipanti posso rivolgersi agli organizzatori per avere le credenziali di accesso
    \item Non è permesso l'utilizzo di \textbf{dispositivi elettronici} diversi dal \textit{PC di gara} e dall'eventuale secondo dispositivo utilizzato come hotspot personale, opportunamente segnalato
    \item Ogni utilizzo del \textit{PC di gara} che non attenga alla risoluzione dei problemi proposti (compresa l'installazione non espressamente autorizzata di software) può comportare \textbf{penalità immediate} fino anche alla squalifica dalla gara
    \item È naturalmente consentito usare carta e penna e \textbf{materiale di cancelleria}
    \item Ogni squadra deve risolvere i \textbf{5 problemi assegnati}. Al termine della gara, le soluzioni di tutti problemi devono essere consegnate alla giuria attraverso un form web (disponibile al momento della gara) rispettando un \textbf{formato ben definito}, specificato nei \textbf{testi dei problemi}
    \item Durante lo svolgimento della gara, ai partecipanti è consentito \textbf{comunicare} solo con i propri compagni di squadra
    \item Alcuni problemi richiedono la \textbf{scrittura di codice}. È consentito utilizzare un semplice editor di testo e un compilatore a riga di comando, oppure avvalersi del servizio web di compilazione e test \href{www.onlinegdb.com}{OnlineGDB}. I \textbf{linguaggi accettati} dalla giuria sono \textsc{C}, \textsc{C++}, \textsc{Javascript} e \textsc{Python}
    \item Gli organizzatori \textbf{vigilano} attivamente per garantire il rispetto delle regole, e possono comminare \textbf{penalità} (fino anche alla \textbf{squalifica} della squadra) nel caso di violazioni del regolamento
    \item Per qualsiasi \textbf{dubbio}, i partecipanti possono rivolgersi \textit{ordinatamente} al team di Street Science
\end{enumerate}


\end{document}
